We consider mixed-integer linear programs of the form
\begin{equation}
\label{eq:milp}
\begin{aligned}
    \min_{x} \quad & c^T x \\
    \text{s.t.} \quad & Ax \leq b, \\
    & l \leq x \leq u, \\
    & x_j \in \mathbb{Z}, \quad \forall j \in \mathcal{I},
\end{aligned}
\end{equation}
where $A \in \mathbb{R}^{m \times n}$, $b \in \mathbb{R}^m$, $c \in \mathbb{R}^n$, and $\mathcal{I} \subseteq \{1, \ldots, n\}$ indexes the variables subject to integrality. The continuous relaxation is obtained by dropping the constraint $x_j \in \mathbb{Z}$ for $j \in \mathcal{I}$ while retaining the bound box $[l, u]$. This relaxation is important because it preserves the large-scale sparse structure of the original MIP and provides a lower bound for minimization problems, but its solutions can lie far from the discrete feasible set.

For the relaxation, we use the standard saddle-point formulation
\begin{equation}
\label{eq:saddle}
    \min_{x \in [l,u]} \max_{y \geq 0} \mathcal{L}(x, y)
    = c^T x + y^T (Ax - b),
\end{equation}
where $y \in \mathbb{R}^m_+$ are the dual multipliers associated with the inequality constraints. The Primal-Dual Hybrid Gradient method updates primal and dual variables through projected first-order steps:
\begin{equation}
\label{eq:pdhg}
\begin{aligned}
    x^{k+1} &= \Pi_{[l,u]}\left(x^k - \tau(c + A^T y^k)\right), \\
    \bar{x}^{k+1} &= x^{k+1} + \theta (x^{k+1} - x^k), \\
    y^{k+1} &= \Pi_{\mathbb{R}^m_+}\left(y^k + \sigma(A\bar{x}^{k+1} - b)\right),
\end{aligned}
\end{equation}
with step sizes $\tau, \sigma > 0$ satisfying $\tau \sigma \|A\|_2^2 < 1$. In the purely convex setting, PDHG enjoys dimension-friendly convergence guarantees and requires only sparse matrix-vector products, which makes it attractive for large-scale optimization \cite{chambolle2011first,lu2023cupdlp}. The difficulty in the present setting is not the relaxation itself, but how to convert its fractional iterates into integer-feasible primal solutions.

\begin{lemma}[Classical PDHG guarantee {\cite{chambolle2011first}}]
\label{lem:pdhg}
Let $(x^k, y^k)$ be generated by \eqref{eq:pdhg} with $\theta = 1$ and $\tau \sigma \|A\|_2^2 < 1$. Then the ergodic averages of the iterates achieve an $\mathcal{O}(1/K)$ primal-dual gap for the continuous relaxation.
\end{lemma}

The lemma explains why PDHG is an effective backbone for the linear relaxation, but it does not address the combinatorial feasibility question that dominates MIP heuristics. To make this gap explicit, we define a joint residual that combines primal infeasibility and non-integrality.

\begin{definition}[Feasibility gap]
\label{def:feasibility-gap}
For $x \in [l,u]$, define the constraint violation vector $v(x) = \max(Ax - b, 0)$ and the integrality residual
\begin{equation}
    \phi(x) = \sum_{j \in \mathcal{I}} \operatorname{dist}(x_j, \mathbb{Z}),
\end{equation}
where $\operatorname{dist}(x_j, \mathbb{Z}) = \min_{k \in \mathbb{Z}} |x_j - k|$. The feasibility gap is
\begin{equation}
    \operatorname{gap}(x) = \max \left\{ \|v(x)\|_\infty, \max_{j \in \mathcal{I}} \operatorname{dist}(x_j, \mathbb{Z}) \right\}.
\end{equation}
\end{definition}

This definition highlights why naive rounding is unreliable: it may reduce $\phi(x)$ while sharply increasing $\|v(x)\|_\infty$. The rest of the paper develops a population-based search procedure that reduces both quantities in a coordinated way. Population dynamics improve coverage of the relaxation landscape, barrier-aware non-local moves remain sensitive to search geometry, and progressive integrality projection reduces the residual without causing an abrupt collapse of feasibility.
